%% Wrote this introduction but I think it's too clunky. Keeping a copy anyway, in case I decide to use it again later.

Elliptic curves are a class of algebraic object long studied in pure mathematics for number-theoretic interest. Since the late 1980s, following the work of Neal Koblitz \cite{koblitz} and others, they have been recognised as a rich source of cryptographic utility; their unique algebraic properties allow small keys to provide a high level of security, and have formed the basis of protocols such as the Elliptic-curve Diffie-Helman which have since obtained widespread adoption. 

These protocols fall under the umbrella of what is called `elliptic-curve cryptography' (ECC). They operate by selecting a suitable elliptic curve and exploiting the arithmetic structure of its points in order to produce cryptographic trapdoor functions. However, these are known to be unsafe against cryptanalytic attacks by quantum computers; for instance, the elliptic-curve discrete logarithm problem---commonly assumed as a source of cryptographic hardness in ECC---is known to be solvable by a quantum computer deploying Shor's algorithm. Therefore, recent interest has pivoted toward more sophisticated cryptographic protocols, still employing the algebraic facilities of elliptic curves, which are conjectured to be quantum-resistant. 

Many proposals to this end exploit the structure of \textit{the space of all elliptic curves} over some field, rather than of just one or another elliptic curve by itself. More specifically, in a suitable context it is possible to consider a graph whose vertices represent isomorphism classes of elliptic curves, and whose edges represent suitable structure-preserving maps between these isomorphism classes. A graph of this kind is called an \textit{isogeny graph}, and has a rich structure, typically consisting of many connected components. One such connected component consists of isomorphism classes of curves having the property of \textit{supersingularity}, which we will later define. This component forms a subgraph which we will later call a \textit{supersingular isogeny graph}.

Supersingular isogeny graphs are a central object of cryptographic interest in the post-quantum view, and are the central object of study in this project. They have desirable properties, such as being regular and being Ramanujan graphs. On this basis, it is conjectured that it is cryptographically hard---even for quantum computers---to find a path joining two arbitrary isomorphism classes of elliptic curves in one of these graphs. Therefore, many proposed post-quantum cryptographic protocols (such as CSIDH, discussed at \cite{CSIDH}) apply what can essentially be understood as random walks on one of these supersingular isogeny graphs. (For more detail, see \cite{de_feo2017}.)

In this project, we will be interested in producing a sample of these supersingular isogeny graphs over a range of parameter values and evaluating properties relevant to their cryptographic applications. In particular, for each such graph we will compute the spectral gap of its adjacency matrix, and mixing times for its random walks. We choose to study these attributes since they are both closely related to the cryptographic hardness conjectured of computing paths on these graphs. 
